% Options for packages loaded elsewhere
\PassOptionsToPackage{unicode}{hyperref}
\PassOptionsToPackage{hyphens}{url}
\documentclass[
  ukrainian,
]{article}
\usepackage{xcolor}
\usepackage[margin=2.5cm]{geometry}
\usepackage{amsmath,amssymb}
\setcounter{secnumdepth}{-\maxdimen} % remove section numbering
\usepackage{iftex}
\ifPDFTeX
  \usepackage[T1]{fontenc}
  \usepackage[utf8]{inputenc}
  \usepackage{textcomp} % provide euro and other symbols
\else % if luatex or xetex
  \usepackage{unicode-math} % this also loads fontspec
  \defaultfontfeatures{Scale=MatchLowercase}
  \defaultfontfeatures[\rmfamily]{Ligatures=TeX,Scale=1}
\fi
\usepackage{lmodern}
\ifPDFTeX\else
\setmainfont{DejaVu Serif}
\setsansfont{DejaVu Sans}
\setmonofont{DejaVu Sans Mono}
\setmathfont{STIX Two Math}
\fi
\ifPDFTeX\else
  % xetex/luatex font selection
\fi
% Use upquote if available, for straight quotes in verbatim environments
\IfFileExists{upquote.sty}{\usepackage{upquote}}{}
\IfFileExists{microtype.sty}{% use microtype if available
  \usepackage[]{microtype}
  \UseMicrotypeSet[protrusion]{basicmath} % disable protrusion for tt fonts
}{}
\makeatletter
\@ifundefined{KOMAClassName}{% if non-KOMA class
  \IfFileExists{parskip.sty}{%
    \usepackage{parskip}
  }{% else
    \setlength{\parindent}{0pt}
    \setlength{\parskip}{6pt plus 2pt minus 1pt}}
}{% if KOMA class
  \KOMAoptions{parskip=half}}
\makeatother
\usepackage{longtable,booktabs,array}
\newcounter{none} % for unnumbered tables
\usepackage{calc} % for calculating minipage widths
% Correct order of tables after \paragraph or \subparagraph
\usepackage{etoolbox}
\makeatletter
\patchcmd\longtable{\par}{\if@noskipsec\mbox{}\fi\par}{}{}
\makeatother
% Allow footnotes in longtable head/foot
\IfFileExists{footnotehyper.sty}{\usepackage{footnotehyper}}{\usepackage{footnote}}
\makesavenoteenv{longtable}
\usepackage{graphicx}
\graphicspath{{/Users/kalinicenkovitalijmikolajovic/.codex/worktrees/5c15/Script_for_translate_book/Output_Final/}}
\setkeys{Gin}{width=0.66\linewidth,height=0.70\textheight,keepaspectratio}
\makeatletter
\newsavebox\pandoc@box
\newcommand*\pandocbounded[1]{% scales image to fit in text height/width
  \sbox\pandoc@box{#1}%
  \Gscale@div\@tempa{\textheight}{\dimexpr\ht\pandoc@box+\dp\pandoc@box\relax}%
  \Gscale@div\@tempb{\linewidth}{\wd\pandoc@box}%
  \ifdim\@tempb\p@<\@tempa\p@\let\@tempa\@tempb\fi% select the smaller of both
  \ifdim\@tempa\p@<\p@\scalebox{\@tempa}{\usebox\pandoc@box}%
  \else\usebox{\pandoc@box}%
  \fi%
}
% Set default figure placement to htbp
\def\fps@figure{htbp}
\makeatother
\ifLuaTeX
\usepackage[bidi=basic,shorthands=off,provide=*]{babel}
\else
\usepackage[bidi=default,shorthands=off,provide=*]{babel}
\fi
\ifLuaTeX
  \usepackage{selnolig} % disable illegal ligatures
\fi
\setlength{\emergencystretch}{3em} % prevent overfull lines
\providecommand{\tightlist}{%
  \setlength{\itemsep}{0pt}\setlength{\parskip}{0pt}}
\usepackage{bookmark}
\IfFileExists{xurl.sty}{\usepackage{xurl}}{} % add URL line breaks if available
\urlstyle{same}
\hypersetup{
  pdflang={uk},
  hidelinks,
  pdfcreator={LaTeX via pandoc}}

\author{}
\date{}

\usepackage{float}
\usepackage[section]{placeins}
\raggedbottom
\clubpenalty=10000
\widowpenalty=10000
\displaywidowpenalty=10000
\brokenpenalty=10000
\emergencystretch=2em
\allowdisplaybreaks[2]
\setlength{\textfloatsep}{12pt plus 2pt minus 4pt}
\setlength{\floatsep}{10pt plus 2pt minus 2pt}
\setlength{\intextsep}{10pt plus 2pt minus 2pt}
\makeatletter
\def\fps@figure{htbp}
\makeatother
\begin{document}

\pandocbounded{\includegraphics[keepaspectratio]{images/_page_205_Figure_5.jpeg}}

Фіг. 28 1. Типова залежність потоку від густини.

Тепер ми готові застосувати \(u_i+f_x=0\) рівняння до проблеми
дорожнього руху. Згідно з правилом диференціації комплексної функції
\(f_x=\frac{df}{du}\,u_x\) ; Це означає, що закон збереження можна
переписати у вигляді

\[
u_t + \frac{df}{du} u_x = u_t + g(u) u_x = 0.
\]

Пусть, например, зависимость потока от плотности квадратичная, т. е.
\(f(u) = u^{2}\) , тогда закон сохранения принимает следующую форму:

\[
u_t + 2uu_x = 0.
\]

Якщо початкова густина розподілу вагонів була \(u(x, 0) = \varphi(x)\) ,
то для знаходження густини розподілу в довільний момент часу необхідно
розв'язати задачу з початковою умовою (задача Коші):

\[
\begin{array}{llll} (\text{УЧП}) & u_t + 2uu_x = 0, & -\infty < x < \infty, & 0 < t < \infty, \\ (\text{HV}) & u\left(x, \ 0\right) = \varphi\left(x\right), & -\infty < x < \infty. \end{array}
\]

Тепер розглянемо розв'язок нелінійної задачі Коші.

\section{Нелінійна задача з початковою умовою (задача
Коші)}\label{ux43dux435ux43bux456ux43dux456ux439ux43dux430-ux437ux430ux434ux430ux447ux430-ux437-ux43fux43eux447ux430ux442ux43aux43eux432ux43eux44e-ux443ux43cux43eux432ux43eux44e-ux437ux430ux434ux430ux447ux430-ux43aux43eux448ux456}

Рассмотрим задачу

\[
\begin{array}{llll} (\mathsf{Y}\mathsf{P}\Pi) & u_t + g\left(u\right)u_x = 0, & -\infty < x < \infty, & 0 < t < \infty, \\ (\mathsf{H}\mathsf{Y}) & u\left(x, \ 0\right) = \varphi\left(x\right), & -\infty < x < \infty. \end{array}
\]

Вспомним предыдущую лекцию, где мы изучали уравнение конвективного
переноса:

\[
u_t + vu_x = 0
\]

.

У цьому рівнянні \(u(x, t)\) --- це концентрація матерії в потоку, що
рухається зі швидкістю v

\[
u_t + g(u)u_x = 0
\]

та рівняння транспорту. Припустимо, що частина води в \(x_0\) рухається
зі швидкістю \(g(u)\) (вздовж потоку або проти нього). Тоді за t секунд
координата частини буде

\[
x = x_0 + g(u) t
\]

(уравнение характеристик).

Пам'ятайте, що концентрація \(u (x, t)\) не змінюється вздовж
характеристики. Якщо ми знаємо початкову концентрацію \(u(x_a, 0)\) ,

\pandocbounded{\includegraphics[keepaspectratio]{images/_page_206_Figure_8.jpeg}}

Фіг. 28.2. Характеристики рівняння \(u_t+g\) (u) \(u_x=0\) ; a ---
загальна форма характеристики, що випливає з характеристики ( \(x_0\) ,
0), вздовж характеристики розв'язок залишається сталим і рівним значенню
\(u (\) x\_0 \(, 0)\); b --- приклад характеристики, що починається на
t. (2, 0), рівняння \(u_t+3uu_x=0\) з \(u (2, 0)\)=1.

то уравнение характеристик принимает вид \(x = x_0 + g[u(x_0, 0)] t\)
(уравнение характеристики, начинающейся в точке \((x_0, 0)\) ).

Пусть, например, мы решаем задачу

(УЧП)

\[
u_t + 3uu_x = 0
\]

,\\
(НУ) u(x, 0) =

\[
\varphi(x)
\]

Щоб знайти характеристику, що виходить із точки, наприклад (2, 0),
потрібно записати

\[
x = 2 + g[u(2, 0)]t =
\]

= \(2 + g[1]t =\) = \(2 + 3t\) .

Можна стверджувати, що розв'язок нашої задачі \(u(x, t)\) дорівнює
одному на лінін x=2+3t, тобто \(u(2+3t, t)\)=1 для всіх t\textgreater0.
Фіг. 28.2 показує характеристики загального випадку та цього конкретного
завдання.

Тепер зрозуміло, що, знаючи характеристики в кожній точці і знаючи, що
розв'язок за характеристиками не змінюється, можна отримати розв'язок
\(u\left(x,\ t\right)\) у будь-який момент часу t. Ми не отримали явного
виразу для \(u\left(x,\ t\right)\) через змінні x і t, але знання
характеристик рівняння дозволяє розв'язати деякі цікаві задачі.

\section{Задача дорожного
движения}\label{ux437ux430ux434ux430ux447ux430-ux434ux43eux440ux43eux436ux43dux43eux433ux43e-ux434ux432ux438ux436ux435ux43dux438ux44f}

Розглянемо модель руху, у якій потік обчислюється за формулою
\(f(u)=u^2\) , і початковою густиною iso-P і C

\pandocbounded{\includegraphics[keepaspectratio]{images/_page_207_Figure_8.jpeg}}

. 28.3. Начальная плотность автомобилей движущихся слева направо.

виведений на Рис. 28.3. Іншими словами, нам потрібно розв'язати проблему

\[
(\text{УЧП}) \quad u_t + 2uu_x = 0, \quad -\infty < x < \infty, \quad 0 < t < \infty,
(\text{HV}) \quad u(x, \ 0) =
\begin{cases} 1, & x \leq 0, \\ 1 - x, & 0 < x < 1, \\ 0, & 1 \leq x. \end{cases}
\]

Почнемо з визначення характеристик, які виникають з \((x_0, 0)\) точок
зору. За \(x_0 < 0\) ми отримуємо

\[
x = x_0 + g [u(x_0, 0)] t =
\]

= \(x_0 + g [1] t =\) = \(x_0 + 2t\) .

Разрешая относительно t, находим уравнение характеристик

\[
t = \frac{1}{2} (x - x_0).
\]

Ці прямі лінії показані на рис. 28.4. Тепер розглянемо точки
\(0 < x_0 < 1\) . Тут характеристики визначаються такими
співвідношеннями:

\[
x = x_0 + g[u(x_0, 0)] t =
= x_0 + g[(1 - x_0)] t =
= x_0 + 2(1 - x_0) t.
\]

Разремая последнее соотношение относительно t, получаем

\[
t = \frac{x - x_0}{2(1 - x_0)}.
\]

При \(1 ≤ x_0 < ∞\) характеристики

\[
x = x_0 + g[u(x_0, 0)]t =
\]

= \(x_0 + g[0]t =\) = \(x_0\) призводить до вертикальних ліній, що
виходять з xo. Усі характеристики задачі показані на рис. 28.4.

\pandocbounded{\includegraphics[keepaspectratio]{images/_page_208_Figure_10.jpeg}}

Фіг. 28.4. Характеристики рівняння \(u_t + 2uu_x = 0\) .

Ситуація на дорозі при 0 \textless{} t \textless{} 1/2 досить зрозуміла.
Рис. 28.5 показує розв'язок нашої задачі в різні моменти часу. Зверніть
увагу, що характеристики сходяться в одній точці t=1/2. Отже, щоб
розв'язати задачу на t\textgreater1/2, потрібно використовувати інший
метод. Коли характеристики сходяться в одній точці, говоримо про ударну
хвилю (розрив розв'язку). Тепер потрібно відповісти на питання: на якій
швидкості є передній край

Ударна хвиля рухатиметься вздовж шосе? Хоча це не очевидно, виявляється,
що швидкість поширення розриву відповідає формулі

\[
S = \frac{f(u_R) - f(u_L)}{u_R - u_L},
\]

де \(u_R\) і \(u_L\) --- значення розв'язку праворуч і ліворуч від
фронту хвилі, а \(f(u_R)\) і \(f(u_L)\) --- значення потоку при цих
густинах.

\pandocbounded{\includegraphics[keepaspectratio]{images/_page_209_Figure_5.jpeg}}

Фіг. 28.5. Щільність руху в різний час.

У нашому прикладі швидкість поширення щілини виявляється рівною одиниці.
Насправді, якщо згадати цей \(f(u) = u^2\) , то одразу отримуємо

\[
S = \frac{0-1}{0-1} = 1
\]

.

Це означає, що при t\textgreater1/2 фронт хвилі рухатиметься зліва
направо зі швидкістю одиниці.

Полное решение задачи

\[
(\text{Y} \, \Pi) \quad u_t + 2uu_x = 0, \quad -\infty < x < \infty,
(\text{НУ}) \quad u(x, \ 0) =
\begin{cases} 1, & x \leq 0, \\ 1 - x, & 0 < x < 1, \\ 0, & 1 < x, \end{cases}
\]

изображено на рис. 28.5. \# ЗАУВАЖЕННЯ 1. Ударна хвиля в нашому прикладі
виникає через те, що потік швидко зростає зі збільшенням густини u. Якщо
потік заданий формулою \(f(u)\) = u, то рівняння було б у вигляді
\(u_1+u_x=0\) , і розв'язок цього рівняння представляє хвилю, що
рухається вправо без спотворень. Якщо на мить замислитися над значенням
формули \(f(u)\) = u, стає очевидно, що розв'язок рухатиметься таким
чином.

\begin{enumerate}
\def\labelenumi{\arabic{enumi}.}
\setcounter{enumi}{1}
\tightlist
\item
  Прямой подстановкой решения в виде
\end{enumerate}

\[
u = \varphi(x - g(u) t)
\]

Ви можете переконатися, що ця формула неявно визначає розв'язок проблеми

\[
\begin{array}{llll} u_t + g\left(u\right) u_x = 0, & -\infty < x < \infty, & 0 < t < \infty, \\ u\left(x, \ 0\right) = \varphi\left(x\right), & -\infty < x < \infty. \end{array}
\]

Например, решение задачи Коши

\[
(Y\Pi\Pi) \quad u_t + uu_x = 0,
(\text{НУ}) \quad u(x, 0) = x
\]

неявно вадается соотношением

\[
u = \varphi [x - g(u) t] =
= x - g(u) t = x - ut.
\]

У нашому конкретному випадку можливо отримати рішення у явній формі (у
багатьох інших випадках це неможливо)

\[
u(x, t) = \frac{x}{1+t}.
\]

Переконайтеся, що ця функція справді є вирішенням нашої проблеми.

\section{\texorpdfstring{\textbf{ЗАВДАННЯ}}{ЗАВДАННЯ}}\label{ux437ux430ux432ux434ux430ux43dux43dux44f}

\begin{enumerate}
\def\labelenumi{\arabic{enumi}.}
\tightlist
\item
  Розв'яжіть задачу Коші
\end{enumerate}

\[
\begin{aligned}
(\text{УЧП})\quad & u_t + uu_x = 0, && -\infty < x < \infty,\quad 0 < t < \infty, \\
(\text{НУ})\quad & u(x,0) =
\begin{cases}
0, & x < 0, \\
x, & 0 \leqslant x.
\end{cases}
\end{aligned}
\]

Побудуйте хронологію рішення для різних моментів часу. Як ви тлумачите
це рішення? Як виглядає зв'язок між потоком і густиною в цій задачі?

\begin{enumerate}
\def\labelenumi{\arabic{enumi}.}
\setcounter{enumi}{1}
\tightlist
\item
  Решите задачу
\end{enumerate}

\begin{verbatim}
\begin{split} & (\text{УЧП}) \quad u_t + u^2 u_x = 0, \quad -\infty < x < \infty, \quad 0 < t < \infty, \\ & (\text{НУ}) \quad u_t(x, 0) = \begin{cases} 0, & x < 0, \\ x, & 0 \leqslant x. \end{cases}
\end{verbatim}

Який зв'язок між потоком і щільністю? Чи співпадає поведінка отриманого
вами рішення з тим, чого ви очікували? Порівняйте рішення задачі 1 і
завдання 2.

\begin{enumerate}
\def\labelenumi{\arabic{enumi}.}
\setcounter{enumi}{2}
\tightlist
\item
  Припустимо, що нев'язка рідина протікає через трубу і одночасно
  просочується крізь стінки згідно із законом \(F(u)\) = ku (г/см·с).
  Закон збереження (або, точніше, закон зміни) у цьому випадку
  записується як
\end{enumerate}

\[
u_t + f_x = -F(u)
\]

.

Знайдіть розв'язок цього рівняння, якщо \(f(u)\) = u і
\(\varphi(x) = u\) (x, 0) (початковий розподіл задається загальним
чином). Дайте тлумачення отриманого рішення.

\begin{enumerate}
\def\labelenumi{\arabic{enumi}.}
\setcounter{enumi}{3}
\item
  Знайти розв'язок попереднього рівняння, якщо втрати визначаються
  функцією \(F(x, t)\) = 1/x. Перевірте розв'язок. Ви розумієте його
  фізичне значення?
\item
  Перевірити, що функція та неявно задана відношенням
\end{enumerate}

\[
u = \varphi[x - g(u) t],
\]

є розв'язком нелінійної задачі

\[
\begin{aligned}
(\text{УЧП})\quad & u_t + g(u)u_x = 0, \\
(\text{НУ})\quad & u(x,0) = \varphi(x).
\end{aligned}
\]

\section{СИСТЕМИ РІВНЯНЬ З ЧАСТКОВИМИ
ПОХІДНИМИ}\label{ux441ux438ux441ux442ux435ux43cux438-ux440ux456ux432ux43dux44fux43dux44c-ux437-ux447ux430ux441ux442ux43aux43eux432ux438ux43cux438-ux43fux43eux445ux456ux434ux43dux438ux43cux438}

Мета лекції - показати, що багато фізичних систем, особливо в
гідродинаміці, не можна описати одним рівнянням. Для цього потрібні
системи пов'язаних рівнянь для тиску, густини, температури та компонент
швидкості.

Лекція покаже, як можна розв'язати лінійну систему рівнянь

\[
u_t + Au_x = 0
\]

якщо перетворити її на нову систему незалежних рівнянь

\[
v_1 + \Lambda v_x = 0,
\]

а затем независимо решить каждое уравнение этой системы.

У багатьох галузях науки існують системи величин, які взаємопов'язані не
одним, а кількома відношеннями одночасно. Наприклад, у гідродинамікі це
чотири рівняння

\[
u_t + uu_x + vu_y + \frac{1}{\rho} p_x = 0
\]

(збереження імпульсу вздовж осі \(x\) ), (29,1)
\(v_t + uv_x + vv_y + \frac{1}{\rho} p_y = 0\) (збереження імпульсу
вздовж осі \(y\) ), (збереження маси), (збереження маси), (збереження
маси),
\(\left(\frac{p}{\rho^{\gamma}}\right)_t + u\left(\frac{p}{\rho^{\gamma}}\right)_x + v\left(\frac{p}{\rho^{\gamma}}\right)_y = 0\)
(збереження енергії), \(p(x, y, z, t)\) --- тиск у рідині,
\(u(x, y, z, t)\) --- \(x\) компонента швидкості, \(v(x, y, z, t)\) ---
\(y\) компонента швидкості, \(\rho(x, y, z, t)\) --- густина рідини.

Задана нелінійна система рівнянь називається рівняннями Ейлера руху
рідини. Завдання полягає в одночасному знаходженні всіх невідомих
функцій p, i, v і p, які задовольняють усім чотирьом рівнянням (разом із
початковими та граничними умовами).

Існують й інші причини вивчати системи рівнянь. У теорії звичайних
диференціальних рівнянь (якщо читач пам'ятає) показано, що одне рівняння
другого порядку можна представити як систему з двох рівнянь першого
порядку. Хоча теорія диференціальних рівнянь у похідних не така проста,
дуже часто одне рівняння з похідними у похідних порядків зводиться до
системи диференціальних рівнянь першого порядку. Наприклад, телеграфне
рівняння

\[
u_{tt} = c^2 u_{xx} + a u_t + b u
\]

путем введения трех новых переменных

\[
u_1 = u_1
u_2 = u_2
u_3 = u_4
\]

можно свести к системе уравнений первого порядка

(29.2)

\[
\begin{cases} \frac{\partial u_1}{\partial x} = u_2, \\ \frac{\partial a_1}{\partial t} = u_3, \\ \frac{\partial u_3}{\partial t} = c^2 \frac{\partial u_2}{\partial x} + au_3 + bu_4. \end{cases}
\]

Для розв'язання систем типу (29.2) необхідно добре знати такі поняття,
як власні значення та власні вектори матриці. З цієї причини ми
пропонуємо короткий огляд лінійної алгебри.

\section{Огляд лінійної
алгебри}\label{ux43eux433ux43bux44fux434-ux43bux456ux43dux456ux439ux43dux43eux457-ux430ux43bux433ux435ux431ux440ux438}

Общий случай

Частный случай

Визначення. Якщо A --- квадратна матриця порядку n, то власні значення
матриці називаються n коренями з

Пусть

\[
A =
\begin{bmatrix} 1 & 1 \\ 4 & 1 \end{bmatrix}
\]

.

ристического уравнения

\[
\det\left(A-\lambda I\right)=0,
\]

где \(\det(A-\lambda I)\) --- определитель матрицы \(A-\lambda I\) .

Деякі власні значення можуть збігатися.

Тогда собственные значения являются корнями уравнения

\[
\det
\begin{bmatrix} 1-\lambda & 1\\ 4 & 1-\lambda \end{bmatrix}
= (1-\lambda)^2 - 4 =
= \lambda^2 - 2\lambda - 3 = 0.
\]

Откуда

\[
\lambda_1 = -1
\]

, \(\lambda_2 = 3\) .

Визначення. Якщо \(\lambda\) --- власне значення матриці A, то
відповідний власний вектор є таким ненульовим вектором, що задовольняє
відношення

\[
Ax = \lambda x
\]

.

Приклад. Вектор X =

\[
\begin{bmatrix} 1 \\ 2 \end{bmatrix}
\]

\(є власним вектором, що відповідає власному значенню\) \lambda\_2 = 3 ,
оскільки

\[
\begin{bmatrix} 1 & 1 \\ 4 & 1 \end{bmatrix}
\begin{bmatrix} 1 \\ 2 \end{bmatrix}
= 3
\begin{bmatrix} 1 \\ 2 \end{bmatrix}
\]

.

Определение. Матрица \(A^{-1}\) называется обратной к матрице A, если

\[
AA^{-1} = A^{-1}A = I_1
\]

Пример, Обратной к

\[
A =
\begin{bmatrix} 2 & -2 \\ 1 & 1 \end{bmatrix}
\]

является матрица

\[
A^{-1} =
\begin{bmatrix} 0.25 & 0.5 \\ -0.25 & 0.5 \end{bmatrix}
\]

,

де I - одинична матриця.

поскольку

\[
AA^{-1} = A^{-1}A =
\begin{bmatrix} 1 & 0 \\ 0 & 1 \end{bmatrix}
\]

.

\section{Діагональна
матриця}\label{ux434ux456ux430ux433ux43eux43dux430ux43bux44cux43dux430-ux43cux430ux442ux440ux438ux446ux44f}

Нехай A --- квадратна матриця порядку n і нехай всі її власні значення
\(\lambda_1\) , \(\lambda_2\) \ldots{} \(\lambda_n\) різні. Побудуємо
матрицю P так, щоб у неї були \emph{стовпці} координати власних векторів
(k-й стовпець містить власний вектор \(X_k\) , що відповідає власному
значенню \(\lambda_k\) ), тобто

\[
P =
\begin{bmatrix} X_1 & X_{\tilde{2}} & \dots & X_n \\ X_{1} & X_{2} & \dots & X_n \end{bmatrix}
=
\begin{bmatrix} x_{11} & x_{1\tilde{2}} & \dots & x_{1n} \\ x_{21} & x_{22} & \dots & x_{2n} \\ \vdots & \vdots & \ddots & \vdots \\ x_{n1} & x_{n\tilde{3}} & \dots & x_{nn} \end{bmatrix}
\]

.

Тогда матрица Л, определяемая соотношением

\[
\Lambda = P^{-1}AP,
\]

где \(P^{-1}\) --- обратная к матрице P, будет диагональной матрицей
вида

\[
\Lambda =
\begin{bmatrix} \lambda_1 & & & \\ & \lambda_2 & 0 & \\ & & \ddots & \\ & & & \lambda_n \end{bmatrix}
\]

.

Например, у матрицы

\[
A =
\begin{bmatrix} 0 & 8 \\ 2 & 0 \end{bmatrix}
\]

власні значення дорівнюють \(\lambda_ 1 = 4\) і \(\lambda_2 = -4\) , а
власні вектори дорівнюють та власним векторам відповідно

\[
X_1 =
\begin{bmatrix} 2 \\ 1 \end{bmatrix}
\]

in X\_2 =

\pandocbounded{\includegraphics[keepaspectratio]{images/_page_221_Picture_9.jpeg}}

Рис. 30.1. Қолеблющаяся мембрана (гиперболическая смешанная задача).

где

\[
\Delta U = U_{rr} + \frac{1}{r} U_r + \frac{1}{r^2} U_{\theta\theta}.
\]

Зверніть увагу, що тут ми вимагаємо, щоб константа розділення була
додатною (тому позначимо її \(\lambda^{2-1}\) )). Лише в цьому випадку
функції \(T(t)\) будуть періодичними.

На следующем шаге мы решим уравнение Гельмгольца, но сначала разберемся
с граничным условием для него. Чтобы найти его, подставим
\(u(r, \theta, t) = U(r, \theta) T(t)\) в граничное условие для мембраны

\[
u(1, \theta, t) = U(1, \theta) T(t) = 0, \quad 0 < t < \infty,
\]

или

\[
U(1, \theta) = 0.
\]

Следовательно, для того чтобы найти форму фундаментальных колебаний
мембраны, необходимо решить задачу

\[
\Delta U + \lambda^2 U = 0,
U(1, \theta) = 0.
\begin{cases} \Delta U + \mu u = 0 & \text{в круге,} \\ U = 0 & \text{на границе круга} \end{cases}
\]

при \(\mu \le 0\) имеет только тривиальное решение U=0. --- Прим. ред.

1) Поскольку краевая задача

Это хорошо известная эллиптическая задача на собственные значения, и нам
надо найти все \(\lambda\) , при которых она имеет ненуленые решения.
Решения \(U\left(r,\,\theta\right)\) описывают форму фундаментальных
колебаний мембраны, а собственные числа \(\lambda\) являются нулями
некоторых функций Бесселя и пропорциональны частотам этих колебаний.

Итак, нам предстоит теперь решить задачу на собственные значения для
уравнения Гельмгольца (эта задача очень важна и сама по себе). Будем
решать ее тем же методом, что и другие линейные однородные уравнения с
нулевыми граничными условиями, а именно методом разделения переменных.

\section{Решение задачи на собственные значения для уравнения
Гельмгольца}\label{ux440ux435ux448ux435ux43dux438ux435-ux437ux430ux434ux430ux447ux438-ux43dux430-ux441ux43eux431ux441ux442ux432ux435ux43dux43dux44bux435-ux437ux43dux430ux447ux435ux43dux438ux44f-ux434ux43bux44f-ux443ux440ux430ux432ux43dux435ux43dux438ux44f-ux433ux435ux43bux44cux43cux433ux43eux43bux44cux446ux430}

Чтобы решить задачу

\[
\begin{array}{llll} (\mathsf{Y}\mathsf{H}\mathsf{\Pi}) & \Delta U + \lambda^2 U = 0, \\ (\mathsf{\Gamma}\mathsf{Y}) & U(1, \theta) = 0, \end{array}
\]

подставим

\[
U(r, \theta) = R(r) \Theta(\theta)
\]

в задачу Гельмгольца. Сделав это, получаем

\[
r^2R'' + rR' + (\lambda^2r^2 - n^2)R = 0
\]

(уравнение Бесселя), \(R(1) = 0\) , \(\Theta'' + n^2\Theta = 0\) (Зроби
все сам). Зверніть увагу, що ми обрали константу розбиття \(n^2\) ,
\(n=0,. \ 1,\ 2,\ \ldots\) тому, що хочемо, щоб функції
\(\Theta\left(\theta\right)\) були періодичними з періодом \(2\pi\) .
Очевидно, що форма мембрани є періодичною функцією \(\theta\) . Отже,
щоб розв'язати задачу Гельмгольца, потрібно розв'язати два звичайних
диференціальні рівняння

\[
r^2R'' + rR' + (\lambda^2r^2 - n^2) R = 0
\]

, \(0 < r < 1\) , \(R(0) < \infty\) (физическое ограничение),
\(R(1) = 0\) , \(\Theta'' + n^2\Theta = 0\) .

\section{Уравнение
Бесселя}\label{ux443ux440ux430ux432ux43dux435ux43dux438ux435-ux431ux435ux441ux441ux435ux43bux44f}

Уравнение

\[
r^2R'' + rR' + (\lambda^2r^2 - n^2)R = 0
\]

жорошо известно в теории обыкновенных дифференциальных уравнений. Это
так называемое уравнение Бесселя. Оно имеет два

линейно-независимых решения \(R_{1}(r) = A J_{n}(\lambda r)\) функция
Бесселя \emph{первого} рода \emph{п}-го порядка,
\(R_{2}(r) = BY_{n}(\lambda r)\) функция Бесселя второго рода n-го
порядка,

и, следовательно, общее решение этого уравнения представимо в виде

\[
R(r) = AJ_n(\lambda r) + BY_n(\lambda r).
\]

Зверніть увагу, що рішення залежать від параметрів n і \(\lambda\) ,
включених у рівняння. Деякі графіки функції Бесселя показані на рис.
30.2. Щоб визначити функції \(J_ n(\lambda r)\) і \(Y_n(\lambda r)\) ,
можливо

\pandocbounded{\includegraphics[keepaspectratio]{images/_page_223_Figure_8.jpeg}}

\pandocbounded{\includegraphics[keepaspectratio]{images/_page_223_Figure_9.jpeg}}

Рис. 30 2 Функции Бесселя. a --- функции Бесселя первого рода;
\(\pmb{\delta}\) --- функции Бесселя второго рода.

використовуйте метод Фробеніуса, який дозволяє знаходити \(R(r)\) у
вигляді силових рядів. Цей метод містить два лінійно незалежні ряди
ступенів \(J_n(\lambda r)\) та \(Y_n(\lambda r)\) . Оскільки функція
\(Y_n(\lambda r)\) не обмежена в r=0, ми обираємо розв'язок у вигляді

\[
R(r) = AJ_n(\lambda r).
\]

Последний шаг в определении \(R\left( r\right)\) связан с использованием
граничного условия \(R\left( \mathbf{l}\right) =0\) для определения всех
\(\lambda\) (в данный

Таблица 30.1

{\def\LTcaptype{none} % do not increment counter
\begin{longtable}[]{@{}lll@{}}
\toprule\noalign{}
Корни & уравнений. & \(J_n(r) = 0\) \\
\midrule\noalign{}
\endhead
\bottomrule\noalign{}
\endlastfoot
------- & ------------ & -------------- \\
\end{longtable}
}

{\def\LTcaptype{none} % do not increment counter
\begin{longtable}[]{@{}lllllllll@{}}
\toprule\noalign{}
Γ & 1 & n & & & & & & \\
\midrule\noalign{}
\endhead
\bottomrule\noalign{}
\endlastfoot
& & 0 & 1 & 2 & 3 & 4 & & \\
& 1 & 2,40 & 3,83 & 5,13 & 6,38 & 7,59 & & \\
ı & 2 & 5,52 & 7,02 & 8,42 & 9,76 & 11,06 & & \\
ľ & 3 & 8,65 & 10,17 & 11,62 & 13,02 & 14,37 & & \\
& 4 & 11,79 & 13,32 & 14,80 & 16,22 & 17,62 & & \\
& 5 & 14,93 & 16,47 & 17,96 & 19,41 & 20,83 & & \\
& . & • & & • & & • & & \\
& & & • & & • & • & & \\
& • & • & • & • & • & • & & \\
\end{longtable}
}

момент нас не интересует коэффициент A). Подставляя \(R (1)\) = 0, в
\(R (r)\) = \(J_n\) ( \(\lambda r\) ), получаем \(J_n\) ( \(\lambda\) )
= 0. Другими словами, для того чтобы функция \(R (r)\) обращалась в нуль
на границе круга, мы должны выбрать константу разделения так, чтобы она
была корнем уравнения \(J_n\) (r) = 0, \(\tau\) . e.

\[
\lambda = k_{nm}
\]

где \(k_{nm}\) --- m-й корень уравнения \(J_n(r) = 0\) .

Таблиці цих коренів добре відомі, але їх можна обчислити на комп'ютері.
Значення перших кількох коренів наведені в Таблиці 1. 30.1. Знаючи ці
корені, ми можемо одразу побудувати розв'язок задачі граничного значення
для рівняння Гельмгольца. Якщо власні значення \(\lambda = k_{nm}\) , то
власні функції \(U_{ nm}(r)\) визначаються формулою

\[
U_{nm}(r, \theta) = J_n(k_{nm}r) [A \sin(n\theta) + B \cos(n\theta)].
\]

На рис. 30.3 изображены некоторые из этих функций для различных значений
n и m. Общая форма функции \(U_{nm}(r,\theta)\) не зависит от значений
констант A и B. Эти константы влияют только на амплитуду колебаний и
расположение линии узлов относительно \(\theta=0\) .

Каждая функция представляет собой фундаментальное колебание круглой
мембраны с частотой

\[
f_{nm} = k_{nm}c/2\pi.
\]

Эти частоты мы нашли из решений

\[
T_{nm}(t) = A \sin(k_{nm}ct) + B \cos(k_{nm}ct)
\]

\section{уравнения}\label{ux443ux440ux430ux432ux43dux435ux43dux438ux44f}

\[
T'' + k_{nm}^2 c^2 T = 0.
\]

Интересно отметить, что отношение частоты колебания
\(U_{nm}(\mathbf{r}, \mathbf{\theta})\) к частоте основного колебания
U\_\{nt\}(\ma

\[
sU(x) = -V \frac{dU}{dx}, \quad 0 < x < \infty,
\begin{cases} u(R_1, \theta) = g_1(\theta), \\ u(R_2, \theta) = g_2(\theta). \end{cases}
U(0) = P/s.
\]

\begin{verbatim}
\begin{cases} 2/3, & n=1, \\0 & \text{для всіх інших } n, \\ d_n=\begin{cases} -2/3, & n=1, \\0 & \text{для всіх інших} n. \end{cases}
\end{verbatim}

\[
H(\xi) =
\begin{cases} 1, & 0 \le \theta < \pi, \\ 0, & \pi \le \theta < 2\pi. \end{cases}
\begin{cases} 0, & \xi < 0, \\ 1, & \xi \geqslant 0. \end{cases}
(\Gamma Y)\quad u(1,\varphi)=\cos(3\varphi).
\]

Используя тригонометрию, попытайтесь выразить \(\cos{(3\phi)}\) через
\(\cos{\phi}, \, \cos^2{\phi}, \, \cos^3{\phi}, \, a\) затем
скомбинируйте их так, чтобы по-

лучилось разложение

\[
\cos(3\varphi) = a_0 P_0(\cos\varphi) + a_1 P_1(\cos\varphi) + \dots
\]

\begin{enumerate}
\def\labelenumi{\arabic{enumi}.}
\setcounter{enumi}{4}
\tightlist
\item
  Решите задачу Дирихле в шаре
\end{enumerate}

\[
\Delta u = 0, \quad 0 < r < 1, \quad 0 \le \varphi \le \pi, \quad 0 \le \theta < 2\pi.
u(1, \varphi) =
\begin{cases} 1, & 0 \le \varphi \le \pi/2, \\ -1, & \pi/2 < \varphi \le \pi. \end{cases}
\]

В этой задаче верхняя полусфера нагрета (+1), а нижняя охлаждена (-1).

\begin{enumerate}
\def\labelenumi{\arabic{enumi}.}
\setcounter{enumi}{5}
\tightlist
\item
  Найдите решение задачи Дирихле в шаре
\end{enumerate}

(УЧП)

\[
\Delta u = 0
\]

, \(1 < r < \infty\) , \(0 \le \varphi \le \pi\) ,
\(0 \le \theta < 2\pi\) , \((\Gamma Y)\)
\(u(1, \varphi) = 1 + \cos \varphi\) .

Проверьте ответ.

Лекция 36

\section{НЕОДНОРОДНАЯ ЗАДАЧА ДИРИХЛЕ (ФУНКЦИЯ
ГРИНА)}\label{ux43dux435ux43eux434ux43dux43eux440ux43eux434ux43dux430ux44f-ux437ux430ux434ux430ux447ux430-ux434ux438ux440ux438ux445ux43bux435-ux444ux443ux43dux43aux446ux438ux44f-ux433ux440ux438ux43dux430}

МЕТА ЛЕКЦІЇ: Показати, як гетерогенну задачу Діріхле можна розв'язати за
допомогою функції Гріна (вихідної функції). У цьому важливому методі
права частина рівняння розглядається як певна вхідна дія і розкладається
на неперервний набір дельтаподібних джерел, розподілених по заданій
площі. Потім визначається відповідь системи на кожне таке джерело
(функція Гріна), і всі відповіді підсумовуються (інтегровані). У
результаті ми отримуємо повне розв'язання проблеми.

Достаточно общей задачей прикладной математики является определение
потенциала в некоторой области пространства, если задано распределение
источников \(f(x, y)\) енутри этой области. В электростатике потенциал
(в вольтах) в области D обусловлен распределением зарядов с плотностью
\(f(x, y)\) по этой области. Типичная задача --- определение потенциала
в круге в двумерном случае.

В такой задаче потенциал должен удовлетворять уравнению

Пуассона (например, с нулевыми граничными условиями)

(36.1) (УЧП)

\[
u_{rr} + \frac{1}{r} u_r + \frac{1}{r^2} u_{\theta\theta} = f(r, \theta), \quad 0 \le r < 1, \quad 0 \le \theta < 2\pi,
\]

(ГУ) \(u(1, \theta) = 0, \quad 0 \le \theta \le 2\pi.\) Ми не випадково
обрали нульові умови. Якщо б ми хотіли розв'язати задачу в загальному
випадку, з ненульовими граничними умовами та неоднорідним рівнянням, то
внесок неоднорідних граничних умов можна було б знайти за допомогою
формули інтегралу Пуассона з лекції 33.

Чтобы получить хоть какое-то представление о неоднородных
дифференциальных уравнениях, давайте рассмотрим графически решение
следующей задачи для уравнения Пуассона:

(УЧП)

\[
\Delta u = -q
\]

, \(0 < r < 1\) , \(0 \le \theta \le 2\pi\) ( \(q\) ---положительная
константа),

(\text{НУ})

\[
u(1, \theta) = 0, \quad 0 \le \theta \le 2\pi.
\]

На межі потенціал (або температура, якщо ви бажаєте) утримується на
нулі, а лапласіан дорівнює q у всіх точках області. Оскільки
\(\Delta u\left(p\right)\) слугує мірою різниці між функцією
\(u\left(p\right)\) і її середнім, рівняння Пуассона стверджує, що
поверхня \(u\left(r\right)\) завжди опукла вниз. Інакше кажучи, вона
виглядатиме як тонка мембрана з нерухомим краєм, яка здувається зверху
вниз потоком повітря.

Якщо права частина рівняння \(f(x, y)\) змінюється від точки до точки,
то опуклість також буде змінною функцією точки цієї області.

Перейдем теперь к главной части нашей лекции: определению функции Грина
и решению уравнения (36.1).

Сначала, однако, мы должны ввести понятие потенциала точеных источников
и стоков.

\section{Потенциалы точечных источников и
стоков}\label{ux43fux43eux442ux435ux43dux446ux438ux430ux43bux44b-ux442ux43eux447ux435ux447ux43dux44bux445-ux438ux441ux442ux43eux447ux43dux438ux43aux43eux432-ux438-ux441ux442ux43eux43aux43eux432}

Щоб розв'язати неоднорідне лінійне рівняння, достатньо розв'язати це
рівняння за допомогою точкового джерела, оскільки розв'язок задачі в
загальному випадку можна знайти, підсумувавши внески з точкових джерел.
Наше завдання --- визначити потенціал, який створює точкове джерело (або
поглинач) у певній області. Ці джерела можна інтерпретувати по-різному.
У теорії теплопровідності джерело можна розглядати як точку, де походить
тепло, а дренаж --- як точку, де поглинає тепло. В електростатиці
джерело можна розглядати як окремий позитивний заряд (протон), а дренаж
--- як окремий негативний заряд (електрон). У будь-якому разі, незалежно
від інтерпретації, ми знайдемо двовимірний потенціал і,

создаваемый единичным точечным источником (поиск трехмерного потенциала
мы оставляем читателю в качестве упражнения).

Припустимо, що єдине точкове джерело +q розташоване в початку Clear that
heat (або щось інше)

будет вытекать из области вдоль радиальных линий, и, следовательно, для
везничны полного потока, протекающего через окружность радиуса r можно
записать

Полный поток, вытекающий через окруж-

HOCTD

\[
= -\int_{0}^{2\pi} u_{r}(r) r d\theta = -2\pi r u_{r}(r).
\]

Но полный поток должен совпадать с количеством тепла, созданного внутри
круга (закон сохранения), т. е.

\[
-2\pi r u_r(r) = q.
\]

Решая это диф реренциальное уравнение относительно функции \(u(r)\),
получаем Р И С

\end{document}
